%==============================================================================
% PARTIDA 04: CONDUCTORES Y CABLES
% Codigos: OE.5.2.3.1 al OE.5.2.3.7
%==============================================================================
\subsection{OE.5.2.3.1 Conductor de Cobre LSOH 1.5 mm2 (Fase + Neutro Alumbrado)}

\subsubsection{Especificaciones Tecnicas}

Conductor de cobre electrolitico, cableado flexible, con aislamiento termoplastico LSOH (Low Smoke Zero Halogen) de 1.5 mm2 de seccion, para tension de servicio 450/750 V. Se utiliza como conductor de fase y neutro en los circuitos de alumbrado C1, C3 y C5. El color del aislante sera rojo o marron para fase y blanco o celeste para neutro.

El conductor LSOH presenta baja emision de humos y ausencia de halogenos en caso de incendio, mejorando la seguridad en instalaciones residenciales.

\subsubsection{Requisitos de Materiales}

\begin{itemize}
  \item Conductor de cobre LSOH 1.5 mm2, color rojo/marron (fase), en rollos de 100 m.
  \item Conductor de cobre LSOH 1.5 mm2, color blanco/celeste (neutro), en rollos de 100 m.
  \item Conectores de empalme tipo "wago" o borneras de conexion para cajas de derivacion.
  \item Cinta aislante vinilica de alta temperatura.
\end{itemize}

\subsubsection{Metodo de Ejecucion}

\begin{enumerate}
  \item Cortar el conductor a la longitud requerida, dejando 0.30 m adicionales en cada extremo para conexiones.
  \item Pasar el conductor a traves de la tuberia PVC de 3/4'' (PV-L) utilizando una guia de alambre o cinta pasacable.
  \item Pelar el aislante en los extremos con pelacables, dejando 1.5 cm de cobre expuesto.
  \item Conectar el conductor en las borneras de las cajas de derivacion o directamente a los bornes de los interruptores y luminarias.
  \item Identificar los conductores con el codigo de colores establecido.
\end{enumerate}

\subsubsection{Control de Calidad}

\begin{itemize}
  \item Verificar la seccion del conductor con calibrador (1.5 mm2).
  \item Comprobar la continuidad del conductor con un multimetro.
  \item Medir la resistencia de aislamiento entre conductores (minimo 1 MOhm a 1000 V).
  \item Inspeccionar que no existan empalmes dentro de las tuberias.
\end{itemize}

\subsubsection{Metodo de Medicion}

Se medira por metro lineal (m) de conductor instalado, incluyendo el tendido, conexiones y pruebas de continuidad.

\subsubsection{Unidad de Medida}

m (metro lineal).

\subsubsection{Forma de Pago}

Pago por metro lineal instalado y verificado mediante prueba de continuidad.

%------------------------------------------------------------------------------
\subsection{OE.5.2.3.2 Conductor de Cobre LSOH 1.5 mm2 (Proteccion a Tierra Alumbrado)}

\subsubsection{Especificaciones Tecnicas}

Conductor de cobre LSOH 1.5 mm2 de color verde o verde-amarillo, utilizado como conductor de proteccion a tierra para los circuitos de alumbrado. Conecta las partes metalicas de las luminarias con la barra de tierra del tablero correspondiente.

El metodo de ejecucion, control de calidad, unidad de medida y forma de pago son identicos a la partida OE.5.2.3.1.

%------------------------------------------------------------------------------
\subsection{OE.5.2.3.3 Conductor de Cobre LSOH 2.5 mm2 (Fase + Neutro Tomacorrientes)}

\subsubsection{Especificaciones Tecnicas}

Conductor de cobre electrolitico, cableado flexible, con aislamiento LSOH de 2.5 mm2 de seccion, para tension 450/750 V. Se utiliza como conductor de fase y neutro en los circuitos de tomacorrientes C2, C4 y C6. La seccion de 2.5 mm2 es la minima exigida por el CNE-U para circuitos de tomacorrientes en instalaciones residenciales.

\subsubsection{Requisitos de Materiales}

\begin{itemize}
  \item Conductor de cobre LSOH 2.5 mm2, color rojo/marron (fase), en rollos de 100 m.
  \item Conductor de cobre LSOH 2.5 mm2, color blanco/celeste (neutro), en rollos de 100 m.
  \item Conectores de empalme tipo "wago" de 2.5 mm2.
  \item Cinta aislante.
\end{itemize}

\subsubsection{Metodo de Ejecucion}

\begin{enumerate}
  \item Cortar, pelar y tender el conductor de 2.5 mm2 a traves de la tuberia PVC pesada (PV-P) de 3/4''.
  \item Conectar fase al borne L del tomacorriente y neutro al borne N.
  \item Asegurar que las conexiones esten firmes y sin deformacion del conductor.
\end{enumerate}

\subsubsection{Control de Calidad}

\begin{itemize}
  \item Verificar la seccion de 2.5 mm2.
  \item Comprobar continuidad y polaridad en cada tomacorriente.
  \item Medir resistencia de aislamiento (minimo 1 MOhm).
\end{itemize}

\subsubsection{Metodo de Medicion}

Se medira por metro lineal (m) de conductor de 2.5 mm2 instalado.

\subsubsection{Unidad de Medida}

m (metro lineal).

\subsubsection{Forma de Pago}

Pago por metro lineal instalado y verificado.

%------------------------------------------------------------------------------
\subsection{OE.5.2.3.4 Conductor de Cobre LSOH 2.5 mm2 (Proteccion a Tierra Tomacorrientes)}

\subsubsection{Especificaciones Tecnicas}

Conductor de proteccion a tierra de 2.5 mm2, color verde o verde-amarillo, para conectar los contactos de tierra de los tomacorrientes a la barra de tierra del tablero. Es obligatorio en todos los tomacorrientes segun el CNE-U y la EM.010.

Metodo de ejecucion, control, medicion y pago identicos a OE.5.2.3.3.

%------------------------------------------------------------------------------
\subsection{OE.5.2.3.5 Conductor de Cobre LSOH 2.5 mm2 (Alimentacion Electrobomba)}

\subsubsection{Especificaciones Tecnicas}

Conductor de cobre LSOH 2.5 mm2 para la alimentacion dedicada de la electrobomba monofasica de 1 HP. Incluye conductores de fase, neutro y proteccion a tierra. El circuito debe ser independiente y protegido con interruptor termomagnetico de 16 A.

Metodo de ejecucion, control, medicion y pago identicos a OE.5.2.3.3.

%------------------------------------------------------------------------------
\subsection{OE.5.2.3.6 Conductor de Cobre LSOH 10 mm2 (Alimentador General)}

\subsubsection{Especificaciones Tecnicas}

Conductor de cobre LSOH 10 mm2 para el alimentador general desde la caja portamedidor hasta el tablero general TG-01. Incluye conductores de fase, neutro y proteccion a tierra. La seccion de 10 mm2 esta dimensionada para soportar la corriente de demanda maxima de 21.73 A con factor de seguridad y reserva para ampliaciones.

\subsubsection{Requisitos de Materiales}

\begin{itemize}
  \item Conductor de cobre LSOH 10 mm2, color negro (fase), en rollos de 50 m.
  \item Conductor de cobre LSOH 10 mm2, color blanco (neutro), en rollos de 50 m.
  \item Terminales de compression tipo ojo para 10 mm2.
  \item Cinta aislante autofundente.
\end{itemize}

\subsubsection{Metodo de Ejecucion}

\begin{enumerate}
  \item Tender los conductores de 10 mm2 a traves de la tuberia PVC pesada (PV-P) de 1''.
  \item Conectar los conductores en la caja portamedidor y en el tablero general TG-01.
  \item Utilizar terminales de compression para las conexiones en los bornes del interruptor general.
  \item Identificar los conductores con cinta de colores (fase: negro/rojo, neutro: blanco).
\end{enumerate}

\subsubsection{Control de Calidad}

\begin{itemize}
  \item Verificar la seccion de 10 mm2 con calibrador.
  \item Medir la resistencia de aislamiento (minimo 1 MOhm).
  \item Comprobar la continuidad de los conductores.
  \item Verificar el torque de apriete de las conexiones (segun fabricante del interruptor).
\end{itemize}

\subsubsection{Metodo de Medicion}

Se medira por metro lineal (m) de conductor de 10 mm2 instalado.

\subsubsection{Unidad de Medida}

m (metro lineal).

\subsubsection{Forma de Pago}

Pago por metro lineal instalado y verificado con pruebas de continuidad y aislamiento.

%------------------------------------------------------------------------------
\subsection{OE.5.2.3.7 Conductor de Cobre Temple Blando 6 mm2 (Conexion Puesta a Tierra)}

\subsubsection{Especificaciones Tecnicas}

Conductor de cobre electrolitico de temple blando, seccion 6 mm2, desnudo (sin aislamiento), para la conexion del sistema de puesta a tierra desde la barra de tierra del tablero general hasta el electrodo de puesta a tierra. El conductor debe ser continuo, sin empalmes, y protegido mecanicamente en los tramos expuestos.

\subsubsection{Requisitos de Materiales}

\begin{itemize}
  \item Conductor de cobre temple blando 6 mm2, desnudo.
  \item Tubo PVC de proteccion mecanica para tramos expuestos.
  \item Conector split-bolt de bronce para conexion al electrodo.
  \item Terminal de compression tipo ojo para conexion a barra de tierra.
\end{itemize}

\subsubsection{Metodo de Ejecucion}

\begin{enumerate}
  \item Tender el conductor de cobre desnudo de 6 mm2 desde la barra de tierra del TG-01 hasta el electrodo de puesta a tierra.
  \item Proteger el conductor con tuberia PVC en los tramos donde pueda estar expuesto a danos mecanicos.
  \item Conectar un extremo a la barra de tierra del tablero mediante terminal de compression.
  \item Conectar el otro extremo al electrodo de puesta a tierra mediante conector split-bolt de bronce.
  \item Aplicar gel conductor en la conexion split-bolt para evitar corrosion.
\end{enumerate}

\subsubsection{Control de Calidad}

\begin{itemize}
  \item Verificar la continuidad del conductor entre la barra de tierra y el electrodo.
  \item Medir la resistencia de puesta a tierra (debe ser menor a 15 Ohms).
  \item Inspeccionar la proteccion mecanica del conductor en tramos expuestos.
\end{itemize}

\subsubsection{Metodo de Medicion}

Se medira por metro lineal (m) de conductor de 6 mm2 instalado.

\subsubsection{Unidad de Medida}

m (metro lineal).

\subsubsection{Forma de Pago}

Pago por metro lineal instalado, incluyendo conectores, proteccion y verificacion de continuidad.
